% sections/08_positioning.tex
% JPHQ-04 — Minimal literature positioning (non-derivative, non-authoritative)

\subsection*{Minimal Literature Positioning}

The results presented in this preprint are positioned as
\emph{structural impossibility theorems}. They are analogous in logical role,
though not in domain, interpretation, or formal language, to No-Go results in
physics, theoretical computer science, and the theory of dynamical systems.

In physics, No-Go theorems delineate regions of formal non-existence under fixed
axioms, symmetries, or invariants, such as forbidden transitions, unattainable
states, or mutually incompatible constraints. In theoretical computer science,
impossibility results identify classes of problems, transformations, or system
properties that cannot be realized under specified computational or structural
assumptions. In dynamical systems theory, analogous results arise from invariant
sets, boundary obstructions, or the absence of admissible trajectories.

Across these domains, the defining logical function of a No-Go theorem is
\emph{negative delimitation}: the precise specification of what cannot occur
without violating the assumed formal framework.

The present work fulfills the same logical function for enterprises modeled as
continua $\KEA$ within the \OC/\kw{ETS.K\_EA.v1.1} formalism. All impossibility
claims are expressed exclusively in terms of OC/ETS primitives—admissible state
sets $\Om$, axes $\Ax$, thresholds $\Th$, cycles $\Cyc$, continuity $\kcont$, and
phase-transition constraints—without importing external semantics, managerial
constructs, optimization criteria, or domain-specific explanatory mechanisms.

The contribution is therefore neither comparative nor empirical. It does not
assert that enterprises behave \emph{like} physical or computational systems,
nor does it transfer laws, models, or causal explanations across domains. The
reference to No-Go reasoning elsewhere is strictly structural: once primitives,
admissibility conditions, and transition constraints are fixed, regions of
formal impossibility follow as internal consequences of the formal system
itself.

Accordingly, the bibliography is intentionally minimal. References are included
solely to situate No-Go reasoning, phase-transition constraints, and structural
invariants as established objects of formal analysis. No cited work is used to
introduce additional primitives, reinterpret OC/ETS concepts, justify the
results empirically, or support claims by analogy or authority.

\begin{remark}
This positioning is deliberately conservative. The validity and significance of
the No-Go theorems stated here depend exclusively on the internal coherence,
explicit assumptions, and falsifiability of the OC/ETS formal system in which
they are formulated, and not on empirical corroboration, cross-domain transfer,
or interpretive analogy.
\end{remark}
